%%=============================================================================
%% Voorwoord
%%=============================================================================

\chapter*{\IfLanguageName{dutch}{Woord vooraf}{Preface}}%
\label{ch:voorwoord}

%% TODO:
%% Het voorwoord is het enige deel van de bachelorproef waar je vanuit je
%% eigen standpunt (``ik-vorm'') mag schrijven. Je kan hier bv. motiveren
%% waarom jij het onderwerp wil bespreken.
%% Vergeet ook niet te bedanken wie je geholpen/gesteund/... heeft
Deze bachelorproef is tot stand gekomen vanuit een persoonlijke ervaring met ADHD en hoe moeilijk het is dingen te plannen. 
Tijdens mijn opleiding Toegepaste Informatica merkte ik dat veel digitale tools ontworpen worden voor een algemene gebruiker,
terwijl bepaalde doelgroepen specifieke noden hebben die vaak over het hoofd worden gezien. Studenten
met ADHD vormen zo'n doelgroep, omdat zij in een academische context vaak extra uitdagingen ervaren bij
planning, organisatie en motivatie.

Het ontwikkelen van een prototype voor een ADHD-vriendelijke planningsapp bood mij de kans om zowel mijn
technische vaardigheden als mijn interesse in user experience en toegankelijk design te combineren.

Graag wil ik mijn promotor, Cem Dhondt, bedanken voor de begeleiding en feedback tijdens het proces.
Daarnaast wil ik mijn co-promotor, Michel Vuijlsteke, bedanken voor de inhoudelijke ondersteuning en
suggesties bij het uitwerken van het concept. Tot slot bedank ik iedereen die mij feedback gaf tijdens het
ontwikkelproces van het prototype.
